\chapter{The Ground State Determines  Dynamics }
\label{gsdd}

In classical mechanics, dynamics is specified by a Hamiltonian function; individual solutions do not determine the Hamiltonian. In quantum mechanics, however, the structure of the Schr\"odinger equation implies a striking inversion: knowledge of a single, sufficiently regular stationary state---specifically the nondegenerate ground state---fixes the potential and thereby all dynamics.

This observation appears in scattered forms across the literature: inverse Sturm--Liouville problems, the Riccati formulation of one-dimensional quantum mechanics, supersymmetric factorizations, and the Hohenberg--Kohn theorem. Yet its conceptual significance is rarely stated explicitly. The aim of this paper is to present a clean, general argument and to clarify its physical and philosophical implications.

Consider a single nonrelativistic particle of mass $m$ on $\mathbb{R}^d$ governed by
\begin{equation}
\hat H = -\frac{\hbar^2}{2m}\nabla^2 + V(\mathbf{x}),
\end{equation}
with $V$ real-valued, locally bounded, and confining enough to admit a normalizable ground state $\psi_0>0$ almost everywhere.

We assume:
\begin{enumerate}
\item The ground state is nondegenerate.
\item $\psi_0(\mathbf{x})$ is twice differentiable and nowhere vanishing.
\item Boundary conditions (e.g. decay at infinity) are known.
\end{enumerate}
These assumptions are standard in both mathematical physics and quantum chemistry.

Let $E_0$ be the ground-state energy. The time-independent Schr"odinger equation reads
\begin{equation}
-\frac{\hbar^2}{2m}\nabla^2\psi_0 + V\psi_0 = E_0\psi_0.
\end{equation}
Solving for the potential yields
\begin{equation}
    \label{pot1d}
V(\mathbf{x}) = E_0 + \frac{\hbar^2}{2m}\frac{\nabla^2\psi_0(\mathbf{x})}{\psi_0(\mathbf{x})}.
\end{equation}
Thus, given $\psi_0$ and $E_0$, the potential is uniquely determined pointwise. Since adding a constant to $V$ merely shifts all energies uniformly, the potential is fixed up to an additive constant even if $E_0$ is unknown.

The positivity of $\psi_0$ ensures that the reconstruction formula is well defined. Any other potential $V'$ admitting the same ground state must coincide with $V$ almost everywhere. Hence the ground state uniquely specifies the Hamiltonian.

Once $V$ is fixed, the entire spectral problem is fixed: all excited eigenstates and eigenvalues are determined, the propagator $U(t)=e^{-i\hat H t/\hbar}$ is fixed, and all expectation values and transition amplitudes follow. Therefore, the ground state determines not merely low-energy physics but the complete time evolution of the system.

In one dimension, the result is a special case of inverse Sturm--Liouville theory, where knowledge of the lowest eigenfunction suffices to reconstruct the potential locally.

Defining the superpotential
\begin{equation}
W(x) = -\frac{\hbar}{\sqrt{2m}}\frac{\psi_0'(x)}{\psi_0(x)},
\end{equation}
we obtain
\begin{equation}
V(x) = W^2(x) - \frac{\hbar}{\sqrt{2m}}W'(x) + E_0,
\end{equation}
showing again that $\psi_0$ encodes the full Hamiltonian structure.

The Hohenberg--Kohn theorem states that the ground-state density uniquely determines the external potential up to a constant. The present result may be viewed as a wavefunction-level strengthening: the full ground-state amplitude determines the Hamiltonian directly, without recourse to variational principles.

Let the ground-state wavefunction be written as
\begin{equation}
\psi_0(\mathbf{x}) = e^{-\phi(\mathbf{x})},
\end{equation}
where $\phi$ is a smooth real-valued function. Positivity of the ground state guarantees that this representation is globally well defined. Substituting into the reconstruction formula yields
\begin{equation}
V(\mathbf{x}) = E_0 + \frac{\hbar^2}{2m}\Big[(\nabla \phi)^2 - \nabla^2 \phi\Big].
\end{equation}
The potential is thus a scalar functional constructed from first and second derivatives of $\phi$, suggesting an intrinsic geometric interpretation.

Define a conformally flat metric on configuration space
\begin{equation}
g_{ij}(\mathbf{x}) = e^{-2\phi(\mathbf{x})},\delta_{ij}.
\end{equation}
The associated Laplace--Beltrami operator satisfies
\begin{equation}
\Delta_g = e^{2\phi}\bigl(\nabla^2 - (d-2)\nabla \phi \cdot \nabla\bigr).
\end{equation}
In this formulation, the ground-state Schr\"odinger equation may be reinterpreted as a geometric condition relating curvature invariants of $g_{ij}$ to the energy $E_0$. The classical potential thus emerges as a curvature scalar of a geometry induced by the quantum state.

Introduce the quantum potential
\begin{equation}
Q[\psi] = -\frac{\hbar^2}{2m}\frac{\nabla^2\psi}{\psi}.
\end{equation}
For the ground state one has
\begin{equation}
V = E_0 - Q[\psi_0].
\end{equation}
Thus the classical potential coincides with minus the quantum potential evaluated on the ground state. From a Hamilton--Jacobi viewpoint, $\phi$ plays the role of an imaginary action, and the ground state defines a preferred foliation of configuration space.

Physical quantum states correspond to rays in Hilbert space, forming a projective manifold equipped with the Fubini--Study metric. The ground state selects a distinguished point in this space. The reconstruction of $V$ from $\psi_0$ may be interpreted as lifting geometric data from projective Hilbert space to a differential-geometric structure on configuration space, reinforcing the view that the Hamiltonian is not primitive but induced by state geometry.

Degenerate ground states may lead to non-uniqueness. Excited states cannot be used directly due to the presence of nodes. For interacting many-body systems the statement generalizes formally, but reconstruction becomes nonlocal in configuration space. The argument does not extend straightforwardly to relativistic or quantum field theoretic settings.

This determinacy challenges the intuition that dynamics is prior to states. In the Schr\"odinger framework, the ground state contains the full dynamical content. The Hamiltonian may therefore be regarded as a derived object, secondary to a privileged stationary state. This perspective aligns naturally with structural and geometric interpretations of quantum theory.


\section{Explicit Example}

Consider the nontrivial ground state
\begin{equation}
    \label{gsx4}
\psi_0(x) = A \exp\Big(-\frac{a}{4} x^4\Big), \quad a>0.
\end{equation}

Substituting \eqref{gsx4} into \eqref{pot1d}, we compute:
\begin{align*}
\psi_0'(x) &= -a x^3 \psi_0(x), \
\psi_0''(x) &= (a^2 x^6 - 3 a x^2) \psi_0(x).
\end{align*}
Thus the potential is
\begin{equation}
V(x) = E_0 + \frac{\hbar^2}{2m} (a^2 x^6 - 3 a x^2),
\end{equation}
a sextic potential with a negative quadratic term, uniquely determined by $\psi_0$.

Writing excited states as $\psi_n = \psi_0 P_n$, the Sturm--Liouville operator for $P_n$ is
\begin{equation}
L = -\frac{\hbar^2}{2m} \left( \frac{d^2}{dx^2} + 2 a x^3 \frac{d}{dx} \right),
\end{equation}
with weight $w(x) = \psi_0^2 = A^2 e^{- \frac{a}{2} x^4}$. The eigenvalue equation is
\begin{equation}
L P_n = (E_n - E_0) P_n,
\end{equation}
which defines the spectrum and orthogonal polynomials (Freud polynomials) relative to $w(x)$.
Numerical or recursive methods can be used to compute $E_n$ and $P_n$.

